% Part: many-valued-logic
% Chapter: three-valued-logics
% Section: kleene

\documentclass[../../../include/open-logic-section]{subfiles}

\begin{document}

\olfileid[hi]{mvl}{thr}{skl}

\olsection{क्लिनी तर्कशास्त्र}

स्टीफन क्लिनी ने ऐसे दृष्टिकोण से प्रेरित दो त्रिमूल्य तर्कशास्त्र प्रस्तुत
किए जिसमें सत्य-मानों को संगणनात्मक प्रक्रियाओं के परिणाम माना जाता है:
कोई प्रक्रिया $\True$ अथवा $\False$ दे सकती है, पर समाप्त होने में विफल भी
हो सकती है। उस स्थिति में अनुरूप सत्य-मान अपरिभाषित है, जिसे सत्य-मान
~$\Undef$ से दर्शाते हैं।

किसी प्रतिज्ञप्ति~$!A$ का निषेध संगणित करने के लिए पहले~$!A$ का मान संगणित
करके परिणाम का विपरीत लौटाएँगे। यदि $!A$ की संगणना समाप्त नहीं होती, तो
पूरी प्रक्रिया भी समाप्त नहीं होती; अतः $\Undef$ का निषेध $\Undef$ है।

संयोजन $!A \land !B$ संगणित करने के दो विकल्प हैं: पहले~$!A$, फिर $!B$
संगणित करें; दोनों का परिणाम~$\True$ होने पर परिणाम $\True$, और अन्यथा
$\False$ होगा। यदि कोई एक संगणना रुकने में विफल हो, तो पूरी प्रक्रिया भी
विफल होती है। अतः इस स्थिति में यदि कोई एक संयोज्य अपरिभाषित है, तो संयोजन
भी अपरिभाषित है। वियोजन के लिए भी यही बात लागू होती है।

किंतु यदि हम $!A$ और $!B$ का मूल्यांकन समांतर रूप से कर सकें, तो बेहतर कर
सकते हैं। तब यदि दोनों प्रक्रियाओं में से कोई रुककर $\False$ लौटाए, तो हम
रुक सकते हैं, क्योंकि उत्तर असत्य ही होगा। अतः उस स्थिति में एक असत्य
संयोज्य वाला संयोजन असत्य है, चाहे दूसरा संयोज्य अपरिभाषित हो। इसी प्रकार,
वियोजन को समांतर संगणित करते समय किसी एक वियोज्य की प्रक्रिया के सत्य
लौटाते ही हम रुक सकते हैं; तब वियोजन सत्य ही होगा। अतः इस स्थिति में कोई
घटक अपरिभाषित होने पर भी हम संयुक्त कथन का परिणाम जान सकते हैं। इस
निर्वचन में $\Undef$ को ``अपरिभाषित'' के बदले ``अज्ञात'' पढ़ सकते हैं।

दोनों निर्वचनों से क्लिनी के प्रबल और दुर्बल तर्कशास्त्र मिलते हैं।
सप्रतिबंध को $\lnot !A \lor !B$ के तुल्य परिभाषित किया जाता है।

\begin{defn}
\emph{प्रबल क्लिनी तर्कशास्त्र}~$\LogKs$ इस आव्यूह से परिभाषित है:
\begin{enumerate}
  \item मानक प्रतिज्ञप्तिपरक भाषा $\Lang L_0$, जिसमें
  $\lnot$, $\land$, $\lor$, $\lif$.
  \item सत्य-मानों का समुच्चय $V = \{\True, \Undef, \False\}$।
  \item $\True$ एकमात्र निर्दिष्ट मान है, अर्थात् $V^+ = \{\True\}$।
  \item सत्य-फलन निम्न सारणियों से दिए जाते हैं:
  \begin{center}
    \begin{tabular}{c|c} 
      $\tf{\lnot}$ & \\ 
      \hline  
      $\True$ & $\False$ \\ 
      $\Undef$ & $\Undef$ \\
      $\False$ & $\True$ 
    \end{tabular}
    \quad
    \begin{tabular}{c|ccc} 
      $\tf{\land}[\LogKs]$ & $\True$ & $\Undef$ & $\False$ \\ 
      \hline 
      $\True$ & $\True$ & $\Undef$ & $\False$ \\ 
      $\Undef$ & $\Undef$ & $\Undef$ & $\False$\\ 
      $\False$ & $\False$ & $\False$ & $\False$ 
    \end{tabular}
    \\[2ex]
    \begin{tabular}{c|ccc} 
      $\tf{\lor}[\LogKs]$ & $\True$ & $\Undef$ & $\False$ \\ 
      \hline 
      $\True$ & $\True$ & $\True$ & $\True$ \\ 
      $\Undef$ & $\True$ & $\Undef$ & $\Undef$ \\
      $\False$ & $\True$ & $\Undef$ & $\False$ 
    \end{tabular}
    \quad
    \begin{tabular}{c|ccc} 
      $\tf{\lif}[\LogKs]$ & $\True$ & $\Undef$ & $\False$ \\ 
      \hline 
      $\True$ & $\True$ & $\Undef$ & $\False$ \\ 
      $\Undef$ & $\True$ & $\Undef$ & $\Undef$  \\ 
      $\False$ & $\True$ & $\True$ & $\True$ 
    \end{tabular}
  \end{center} 
\end{enumerate}
\end{defn}

\begin{defn}
\emph{दुर्बल क्लिनी तर्कशास्त्र}~$\LogKw$ इस आव्यूह से परिभाषित है:
\begin{enumerate}
  \item मानक प्रतिज्ञप्तिपरक भाषा $\Lang L_0$, जिसमें
  $\lnot$, $\land$, $\lor$, $\lif$.
  \item सत्य-मानों का समुच्चय $V = \{\True, \Undef, \False\}$।
  \item $\True$ एकमात्र निर्दिष्ट मान है, अर्थात् $V^+ = \{\True\}$।
  \item सत्य-फलन निम्न सारणियों से दिए जाते हैं:
  \begin{center}
    \begin{tabular}{c|c} 
      $\tf{\lnot}$ & \\ 
      \hline  
      $\True$ & $\False$ \\ 
      $\Undef$ & $\Undef$ \\
      $\False$ & $\True$ 
    \end{tabular}
    \quad
    \begin{tabular}{c|ccc} 
      $\tf{\land}[\LogKw]$ & $\True$ & $\Undef$ & $\False$ \\ 
      \hline 
      $\True$ & $\True$ & $\Undef$ & $\False$ \\ 
      $\Undef$ & $\Undef$ & $\Undef$ & $\Undef$\\ 
      $\False$ & $\False$ & $\Undef$ & $\False$ 
    \end{tabular}
    \\[2ex]
    \begin{tabular}{c|ccc} 
      $\tf{\lor}[\LogKw]$ & $\True$ & $\Undef$ & $\False$ \\ 
      \hline 
      $\True$ & $\True$ & $\Undef$ & $\True$ \\ 
      $\Undef$ & $\Undef$ & $\Undef$ & $\Undef$ \\
      $\False$ & $\True$ & $\Undef$ & $\False$ 
    \end{tabular}
    \quad
    \begin{tabular}{c|ccc} 
      $\tf{\lif}[\LogKw]$ & $\True$ & $\Undef$ & $\False$ \\ 
      \hline 
      $\True$ & $\True$ & $\Undef$ & $\False$ \\ 
      $\Undef$ & $\Undef$ & $\Undef$ & $\Undef$  \\ 
      $\False$ & $\True$ & $\Undef$ & $\True$ 
    \end{tabular}
  \end{center} 
\end{enumerate}
\end{defn}

\begin{prop}
  $\LogKs$ और $\LogKw$ की कोई सर्वसत्यता नहीं है।
\end{prop}

\begin{proof}
  यदि सभी !!{propositional variable}~$p$ के लिए $\pAssign v(p) =
  \Undef$, तो किसी भी सूत्र~$!A$ का सत्य-मान~$\pValue v(!A) = \Undef$
  होगा, क्योंकि
  \[
    \tf{\lnot}(\Undef) = \tf{\lor}(\Undef, \Undef) = \tf{\land}(\Undef,
  \Undef) = \tf{\lif}(\Undef, \Undef) = \Undef
  \]
  दोनों तर्कशास्त्रों में। चूँकि $\LogKs$ और $\LogKw$ दोनों के लिए $\Undef
  \notin V^+$, इसलिए इस !!{valuation} में $!A$ निर्दिष्ट नहीं होगा।
\end{proof}

यद्यपि दुर्बल और प्रबल दोनों क्लिनी तर्कशास्त्रों की कोई सर्वसत्यता नहीं,
उनके परिणाम-संबंध अतुच्छ हैं।

\begin{prob}
  निम्न में से कौन-से संबंध (क) प्रबल और (ख) दुर्बल क्लिनी तर्कशास्त्र में
  लागू होते हैं? प्रत्येक के लिए सत्य-सारणी दें।
  \begin{enumerate}
    \item $p, p \lif q \Entails q$
    \item $p \lor q, \lnot p \Entails q$
    \item $p \land q \Entails p$
    \item $p \Entails p \land p$
    \item $p \Entails p \lor q$
  \end{enumerate}
\end{prob}

दिमित्री बोचवार ने $\Undef$ का निर्वचन ``अर्थहीन'' के रूप में किया और यह
नियम रखकर मिथ्यावादी विरोधाभास जैसे विरोधाभासों को हल करने का प्रयास किया
कि विरोधाभासी वाक्य~$\Undef$ मान लें। उन्होंने ऐसा तर्कशास्त्र प्रस्तुत
किया जो मूलतः अतिरिक्त संयोजकों से विस्तारित दुर्बल क्लिनी तर्कशास्त्र है;
इनमें दो हैं ``बाह्य निषेध'' और ``अपरिभाषित है'' संक्रियक:
\begin{center}
  \begin{tabular}{c|c} 
    $\tf{\sim}$ & \\ 
    \hline  
    $\True$ & $\False$ \\ 
    $\Undef$ & $\True$ \\
    $\False$ & $\True$ 
  \end{tabular}
\quad
  \begin{tabular}{c|c} 
    $\tf{+}$ & \\ 
    \hline  
    $\True$ & $\False$ \\ 
    $\Undef$ & $\True$ \\
    $\False$ & $\False$ 
  \end{tabular}
\end{center}

\begin{prob}
  क्या आप बोचवार के तर्कशास्त्र में $\sim$ को $\lnot$ और~$+$ के पदों में
  परिभाषित कर सकते हैं? अर्थात् केवल !!{propositional variable}~$p$ वाला,
  $\sim$ रहित ऐसा !!a{formula} खोजें जो सदा~$\mathord{\sim}p$ वाला ही
  सत्य-मान ले। अपनी बात दिखाने के लिए सत्य-सारणी दें।
\end{prob}


\end{document}
